\par\Needspace{6\baselineskip}\medskip
\hypertarget{native-211-article156}{}\pdfbookmark[2]{156.}{article-156}
\begin{center}156.\end{center}
Exemplum \textit{secundum} a Pallade sumemus, cuius observationes sequentes Mediolani factas e Commercio litterario clar.\ \textsc{de Zach}, Vol.\ XIV. pag.\ 90 excerpimus.
\begin{center}
\begingroup\footnotesize\setlength{\tabcolsep}{4pt}\renewcommand{\arraystretch}{1.2}
\begin{tabular}{|l|r|l|}\hline
Tempus medium Mediol. & Asc.\ recta app. & Declin.\ app.\\\hline
1805 Nov.\ 5.\ $14^{\mathrm h}\,14^{\mathrm m}\,4^{\mathrm s}$ & $78^{0}\,20'\,37''{,}8$ & $27^{0}\,16'\,56''{,}7$ Austr.\\
\hspace{2.4em}Dec.\ 6.\ $11\quad51\quad27$ & $73\quad8\quad48{,}8$ & $32\quad52\quad44{,}3$\\
1806 Ian.\ 15.\ $8\quad50\quad36$ & $67\quad14\quad11{,}1$ & $28\quad38\quad8{,}1$\\\hline
\end{tabular}\endgroup
\end{center}

Loco eclipticae hic aequatorem tamquam planum fundamentale accipiemus, calculoque ita defungemur, ac si orbita penitus adhuc incognita esset. Primo e tabulis Solis pro temporibus propositis sequentia petimus:
\begin{center}
\begingroup\footnotesize\setlength{\tabcolsep}{4pt}\renewcommand{\arraystretch}{1.2}
\begin{tabular}{|l|r|r|r|}\hline
& \shortstack{Longitudo Solis ab\\aequin.\ med.} & Distantia a terra & \shortstack{Latitudo\\Solis}\\\hline
Nov.\ 5 & $223^{0}\,14'\,7''{,}61$ & $0{,}990\,4311$ & $+0''{,}59$\\
Dec.\ 6 & $254\quad28\quad42{,}59$ & $0{,}984\,6753$ & $+0{,}12$\\
Ian.\ 15 & $295\quad5\quad47{,}62$ & $0{,}983\,8153$ & $-0{,}19$\\\hline
\end{tabular}\endgroup
\end{center}

Longitudines Solis, adiectis praecessionibus $+7''{,}59$, $+3''{,}36$, $-2''{,}11$ ad initium anni 1806 reducimus, ac dein, adhibita obliquitate media $23^{0}\,27'\,53''{,}53$ latitudinumque ratione rite habita, ascensiones rectas et declinationes inde deducimus. Ita invenimus
\SourceLeaf{212}
\begin{center}
\begingroup\small\setlength{\tabcolsep}{5pt}\renewcommand{\arraystretch}{1.2}
\begin{tabular}{|l|r|l|}\hline
& Ascensio recta Solis & Declinatio Solis\\\hline
Nov.\ 5 & $220^{0}\,46'\,44''{,}65$ & $15^{0}\,49'\,43''{,}94$ Austr.\\
Dec.\ 6 & $253\quad9\quad23{,}26$ & $22\quad33\quad39{,}45$\\
Ian.\ 15 & $297\quad2\quad51{,}11$ & $21\quad8\quad12{,}98$\\\hline
\end{tabular}\endgroup
\end{center}

Hae positiones ad centrum terrae referuntur, adeoque parallaxi adiecta ad locum observationis reducendae sunt, quum positiones planetae a parallaxi purgare non liceat. Rectascensiones ipsius zenith in hoc calculo adhibendae cum rectascensionibus planetae conveniunt (quoniam observationes in ipso meridiano sunt institutae), declinatio vero ubique erit altitudo poli $=45^{0}\,28'$. Hinc eruuntur numeri sequentes:
\begin{center}
\begingroup\footnotesize\setlength{\tabcolsep}{3pt}\renewcommand{\arraystretch}{1.2}
\begin{tabular}{|l|r|l|r|}\hline
& Asc.\ recta terrae & Declinatio terrae & Log.\ dist.\ a Sole\\\hline
Nov.\ 5 & $40^{0}\,46'\,48''{,}51$ & $15^{0}\,49'\,48''{,}59$ Bor. & $9{,}995\,8375$\\
Dec.\ 6 & $73\quad9\quad23{,}26$ & $22\quad33\quad42{,}83$ & $9{,}993\,3099$\\
Ian.\ 15 & $117\quad2\quad46{,}09$ & $21\quad8\quad17{,}29$ & $9{,}992\,9259$\\\hline
\end{tabular}\endgroup
\end{center}

Loca observata Palladis a nutatione et aberratione fixarum liberanda, ac dein adiecta praecessione ad initium anni 1806 reducenda sunt. Hisce titulis sequentes correctiones positionibus observatis applicare oportebit:
\begin{center}
\begingroup\footnotesize\setlength{\tabcolsep}{4pt}\renewcommand{\arraystretch}{1.2}
\begin{tabular}{|l|r|r||r|r||r|r|}\hline
& \multicolumn{2}{c||}{Observatio I} & \multicolumn{2}{c||}{Observatio II} & \multicolumn{2}{c|}{Observatio III}\\\hline
& Asc.\ r. & Decl. & Asc.\ r. & Decl. & Asc.\ r. & Decl.\\\hline
Nutatio & $-12''{,}86$ & $-3''{,}08$ & $-12''{,}21$ & $-3''{,}42$ & $-13''{,}06$ & $-3''{,}75$\\
aberratio & $-18{,}13$ & $-9{,}89$ & $-22{,}68$ & $-1{,}63$ & $-15{,}60$ & $+9{,}76$\\
praecessio & $+5{,}43$ & $+0{,}62$ & $+2{,}25$ & $+0{,}39$ & $-1{,}51$ & $-0{,}33$\\\hline
summa & $-25{,}56$ & $-12{,}35$ & $-32{,}64$ & $-4{,}66$ & $-30{,}17$ & $+5{,}68$\\\hline
\end{tabular}\endgroup
\end{center}

Hinc prodeunt positiones sequentes Palladis, calculo substruendae:
\begin{center}
\begingroup\small\setlength{\tabcolsep}{6pt}\renewcommand{\arraystretch}{1.2}
\begin{tabular}{|r|r|r|}\hline
T.\ m.\ Parisinum & Asc.\ recta & Declinatio\\\hline
Nov.\ $5{,}574\,074$ & $78^{0}\,20'\,12''{,}24$ & $-27^{0}\,17'\,9''{,}05$\\
$36{,}475\,035$ & $73\quad8\quad16{,}16$ & $-32\quad52\quad48{,}96$\\
$76{,}349\,444$ & $67\quad13\quad40{,}93$ & $-28\quad38\quad2{,}42$\\\hline
\end{tabular}\endgroup
\end{center}
